% RAW dependence resolved by forwarding from EX/MEM to EX.
\begin{tikzpicture}[x=0.95cm,y=1.25cm, font=\footnotesize\sffamily,
  st/.style={draw, minimum width=0.9cm, minimum height=0.5cm, inner sep=0pt, fill=ln-box},
  >={Stealth[length=1.6mm]}]
  \foreach \c in {1,...,6} { \node at (\c,0.7) {\c}; }
  \node[anchor=east] at (0.5,0) {\texttt{ADD R1,R2,R3}};
  \foreach \s/\lab in {0/IF,1/ID,2/EX,3/MEM,4/WB} { \node[st] (a\s) at (\s+1,0) {\lab}; }
  \node[anchor=east] at (0.5,-1) {\texttt{SUB R4,R1,R5}};
  \foreach \s/\lab in {0/IF,1/ID,2/EX,3/MEM,4/WB} { \node[st] (b\s) at (\s+2,-1) {\lab}; }
  %% The result of ADD leaves EX (a2) and is fed into the EX of SUB (b2).
  \draw[->,thick,ln-tip] (a2.south) to[out=-90,in=90,looseness=0.9] (b2.north);
  %% The label sits in the gap between the rows, right of the arrow.  y=1.25
  %% keeps the gap (0.75cm) wide enough for the arrow and the wider Japanese
  %% word, which would otherwise print over the boxes of the SUB row.
  \node[ln-tip, font=\scriptsize, anchor=west] at (4.2,-0.5)
    {\iflangja{フォワーディング}{forwarding}};
\end{tikzpicture}
